\section{前提の関門と自分側の補足}\label{sup:self}
前提として確かめた自己位置の追跡は，経験50\%条件では6組すべてで基準（決定係数 0.80）を満たした．
経験80\%条件は6組とも基準に届かなかった．
共通テスト（ランダムに歩くデータ）では 0.70〜0.75（本番の3組では 0.72〜0.75）であったが，自分の学習データと同じ作りのテストでは 0.99 であり，角にとどまる学習データの偏りによる，経験していない場面への一般化の弱さと考えられる．
% 出典：v13 まとめ §3（G1 の不足）・§10 表（G1 80% 0/6，50% 6/6）。
事前登録した手続きでは，この前提は段階IIへ進む条件であり，経験80\%条件の段階IIは，この前提と G5a の関門を満たさなかったため探索として行った（\ref{app:component}節）．
自分側の積み上がり，G5b，符号に従う分，照合の指標は内部状態からの位置の読み出しを使わず，モデル自身の予測から測るので，追跡が基準に届かないことはこれらの指標の計算を妨げない．
ただし，経験80\%条件のモデルは共通テストの場面に十分には一般化していない状態にあり，以下の経験80\%条件の値は，その状態を含んだものとして読む必要がある．
運動が隠れたステップで予測が赤の方へ一歩動いた場合の視界に寄るか（\ref{app:gates}節）は，経験50\%条件・経験80\%条件とも6組すべてで成り立った．
% 出典：v13 まとめ §10 (a) G2。

自分側の積み上がり $g$ は，角から離れた場面だけを取り出しても，同じ向きで同程度の大きさであり，角の近くだけで働く形ではなかった．
% 出典：v13 まとめ §3（角の近くかどうかで分けた記述）。

ただし，規則の割合を上げると，赤から距離5以内にいるフレームの割合も約0.03，約0.33，約0.85と大きく変わる（本文3.1節）．
角から離れた場面でも同じ向きの結果が得られたことは，角の近くだけで働く形ではないことを示すが，学習時の分布の違いによる影響を除くものではない．
そこで，同じ軌跡で価値と行動の結びつきだけを壊した対照を作った（本文4.3節）．
その対照では，自分側の見込みの快・不快への依存は消えたが，快・不快によらず赤へ向かう癖は一部残った．
また，経験の量は3水準だけであり，強まり方が比例的か飽和的かは分からない．
